\chapter{Introduction and Overview}
\label{chap:intro}

\section{Purpose of this Handbook}

This handbook documents \textbf{piecewise potential-vorticity inversion}
(PPVI) for atmospheric blocking, as implemented in
\texttt{pv\_inversion/}.  The goal is to take the 2025-01-08 California
blocking event, partition its potential-vorticity (PV) anomaly into
vertical pieces, invert each piece to its own balanced circulation, and
thereby quantify how the upper-tropospheric PV anomaly, the
mid/low-tropospheric PV, and the \emph{surface thermal} anomaly each
contribute to the blocking flow.

The pipeline has a single numerical core: \textbf{W.~Wu's
finite-difference successive-over-relaxation (SOR) solver}
(\texttt{wu/fortran/pvpialln\_94UV.f}, \texttt{qinvert21\_94.f},
\texttt{qinvertp21\_94.f}), driven from a thin Python orchestration
layer (\texttt{wu/steps/}).  This Fortran is the same lineage as
C.~Davis's \texttt{PVPI.FOR} \cite{Davis1992} and the modern GFDL/CSU
PPVI tool; the inversion mathematics is unchanged, only the driving and
post-processing are in Python.

\begin{keybox}[What this handbook is about]
A \emph{single} working track --- the Wu Fortran pipeline --- on a
\textbf{nine-level} ERA5 grid, with the PV perturbation split into
\textbf{three pieces} that mirror the reference partition
$\{1\},\{2,3\},\{4..N_L\}$:
\begin{center}\small
\begin{tabular}{lll}
\toprule
piece & levels $K$ & physical content \\
\midrule
\textbf{surface} & $\{1\}$              & 1000\,hPa boundary $\theta$ (the surface heating) \\
\textbf{lower}   & $\{2,3\}$            & 850, 700\,hPa interior PV \\
\textbf{upper}   & $\{4,5,6,7,8,9\}$    & 500--200\,hPa interior PV $+$ 100\,hPa top $\theta$ \\
\bottomrule
\end{tabular}
\end{center}
The key conceptual point (\Cref{chap:math,chap:wu_solver}): the
\textbf{surface warm-$\theta$ anomaly is a boundary condition, not an
interior PV value} --- by Bretherton's equivalence it acts as a
PV \emph{sheet} at the lower boundary, and we invert it as its own
``surface'' piece.
\end{keybox}

\begin{warnbox}[Retired material]
Earlier drafts described a second ``Part~II'' spherical-harmonic (SH)
Python port and several Family-B solvers (a direct spectral solve and an
LGMRES/linearized-Charney solver).  \textbf{These have been removed.}
The SH BALNC/BALP port never converged and is no longer maintained;
this handbook documents only the Wu Fortran pipeline that runs
end-to-end.
\end{warnbox}

\section{The Case: 2025-01-08 California Block}

The event is a strong wintertime ridge over the northeast Pacific / west
coast of North America that contributed to the January 2025 California
wildfire outbreak.  At the analysis time (2025-01-08 00\,UTC) the
500\,hPa block centre is at $\approx 31.5^\circ$N, $222^\circ$E
($138^\circ$W), with $Z_{500}\approx 5890$\,m.  All cross-sections and
the 3-D structure figure (\Cref{sec:step11}) are centred on this point.

\section{Notation Cheat-Sheet}

\begin{center}
\begin{tabular}{ll}
  \toprule
  Symbol & Meaning \\
  \midrule
  $\psi$, $\Phi$ & streamfunction (m$^2$\,s$^{-1}$), geopotential (m$^2$\,s$^{-2}$) \\
  $\bpsi$, $\bPhi$ & climatological mean (basic) state \\
  $\dpsi=\psi-\bpsi$, $\dPhi=\Phi-\bPhi$ & anomalies \\
  $q$ & Ertel potential vorticity (PVU = $10^{-6}\,$K\,m$^2$\,kg$^{-1}$\,s$^{-1}$) \\
  $\bq$, $\dq=q-\bq$ & mean and anomalous PV \\
  $\theta_b$ & boundary potential temperature (surface / top) \\
  $\zeta=\Lap\psi$ & relative vorticity \\
  $f=2\Omega\sin\phi$, $f_0=f(45^\circ)$ & Coriolis parameter, reference value \\
  $\pi=(p/p_0)^{\kappa}$ & Exner function ($p_0=10^{5}$ Pa, $\kappa=R/c_p$) \\
  $\theta=T(p_0/p)^{\kappa}$ & potential temperature \\
  $K=1\ldots N_L$ & 1-indexed Fortran level ($K{=}1$ surface, $K{=}N_L$ top) \\
  $\Lap$, $\Lapinv$ & horizontal Laplacian and its inverse \\
  \bottomrule
\end{tabular}
\end{center}

A more complete glossary appears in \Cref{app:notation}.

\section{Quick Tour of the Code}

\begin{center}\small
\begin{tabular}{lll}
  \toprule
  Component & Purpose & Reference \\
  \midrule
  \texttt{config.py} / \texttt{wu/wu\_config.yaml} & grid, 9 levels, pieces, solver params & \Cref{chap:wu_solver} \\
  \texttt{wu/fortran/pvpialln\_94UV.f}  & Pass A/B: Ertel PV, $\bar\psi$, $\theta_b$ & \Cref{chap:wu_solver} \\
  \texttt{wu/fortran/qinvert21\_94.f}   & Pass C: total balanced inversion (BALNC) & \Cref{chap:wu_solver} \\
  \texttt{wu/fortran/qinvertp21\_94.f}  & Pass D: piecewise perturbation (BALP)    & \Cref{chap:wu_solver} \\
  \texttt{wu/steps/05\ldots11/}         & the step-by-step pipeline + plots         & \Cref{chap:wu_pipeline} \\
  \texttt{wu/steps/08/parse\_and\_pv.py} & Wu\,$Q\to$ SI Ertel PV ($\times10^{-8}$)  & \Cref{chap:wu_coeffs} \\
  \bottomrule
\end{tabular}
\end{center}

\noindent The shared front end (\texttt{shared\_steps/01\ldots04}) downloads
the ERA5 event hour, builds a 30-year climatology, and writes the
Fortran \texttt{.grid} input files; the Wu steps (05--11) compile and run
the four passes and produce all diagnostics.
